\section{Reproducing the measurements}
\label{app:repro}

All commands run from a checkout of the repository \citep{runline2026repo}
with its pinned environment
(\texttt{pip install -r requirements.txt -e '.[dev,build]'}). Results are
written to the git-ignored \texttt{output/} directory.

\begin{small}
\begin{verbatim}
python -m pytest -q -rs                                # Python suite (223 tests)
node --test tests/web_flow.mjs tests/run_session.mjs   # browser logic (35 tests)
python scripts/acceptance.py --data <dir with graphs/ and pools/>
                                    # US1-US3, US5-US7, US10 (local); seed 4971
python scripts/shape_audit.py                          # roundness, dead ends, repeats
python scripts/audit_crossings.py                      # independent crossing labels
python scripts/compare_pools.py                        # 600-request pool check, seed 7
python scripts/validate_deployed.py --pause 1.0 --budget 250
                                    # US4, US9, US10, US11 against production
node scripts/ui_display_check.mjs output/deployed-validation   # US4/US9 UI checks
python scripts/presentation_check.py --api <url> --require-pool
\end{verbatim}
\end{small}

\noindent The acceptance results were produced on code \texttt{61f45cd}
with the rebuilt Charlottesville graph (SHA-256 prefix
\texttt{66b524897627167e}) and its matching complete pool file. The
deployed-system run was on 16 November against backend commit
\texttt{ebcd150}. Full tables, definitions and per-case outputs are in
\texttt{docs/acceptance.md} and \texttt{docs/deployed-validation.md} in the
repository.

\section{Manual checklists}
\label{app:checklists}

\subsection{US8 browser checklist}
Each member runs this once on a desktop browser and once on a phone, at
\url{https://runline-nine.vercel.app}, using only public starts. Record P
(pass), F (fail, with a note) or N/A.

\begin{table}[ht]
\centering\small
\begin{tabularx}{\linewidth}{@{}r Y c c c@{}}
\toprule
\# & Step & Adithya & Krish & Rohan \\
\midrule
1 & Search for ``Rotunda'' and pick the suggestion; the pin moves there. & & & \\
2 & Click elsewhere on the map; the start moves to the click. & & & \\
3 & Request 5 miles, mixed, balanced; routes appear without an error. & & & \\
4 & Switch between the route cards; the map highlights the chosen route. & & & \\
5 & Change surface to trail and regenerate; the results change or a notice explains why not. & & & \\
6 & Change elevation to hilly and regenerate; the climb shown changes or a notice explains why not. & & & \\
7 & Download GPX for the selected route; a file is saved. & & & \\
8 & Enter an invalid distance (0 or blank); an explained error appears. & & & \\
9 & Choose a start outside the supported areas; an explained error appears. & & & \\
10 & Use Start $\rightarrow$ finish with two points in Charlottesville; a one-way route appears. & & & \\
\bottomrule
\end{tabularx}
\end{table}

\subsection{US7 GPX import}
For one route at each of 3, 5, 8 and 12 miles, import the downloaded file into
an outside application \todo{name the application, e.g. Garmin Connect,
Strava or gpx.studio} and record whether the track appears, whether its
distance matches RunLine's within 2\%, and whether start and finish meet.

\subsection{Route review}
Each member independently marks twelve benchmark routes (three starts
$\times$ four distances, routes listed in the acceptance output) on five
questions: (1) the route is continuous; (2) it has no obvious backtracking;
(3) it has no crossing a runner would refuse; (4) it renders correctly on the
map; (5) the displayed metrics agree with what the map shows. Disagreements
are discussed and the agreed outcome recorded.

\section{Detailed results}
\label{app:results}

\paragraph{US11 per area.} Routes within 5\% of target out of 36 (12 requests
at 3, 5 and 8 miles, three routes each). All 216 were closed, within the
$\pm$0.25 mile band and within 10\%, with matching GPX.

\begin{center}\small
\begin{tabular}{@{}l r l@{}}
\toprule
Area & Within 5\% & Answered from \\
\midrule
Washington, DC & 36/36 & precomputed \\
New York & 35/36 & 4 of 12 precomputed, 8 live \\
Los Angeles & 36/36 & precomputed \\
Atlanta & 34/36 & precomputed \\
Boston & 35/36 & precomputed \\
Phoenix & 34/36 & precomputed \\
\bottomrule
\end{tabular}
\end{center}

\paragraph{US6 ties (live).} The eleven pairs where flat and hilly returned
the same top route: ring\_060 5\,mi; ring\_120 3 and 12\,mi; ring\_180 3, 5, 8
and 12\,mi; ring\_240 5, 8 and 12\,mi; ring\_300 3\,mi. No pair had the hilly
route climbing less. After the 18 November fix (\#38) no pair tied; in three
the flat route climbed more (ring\_120 3\,mi, 355 vs 231\,ft; ring\_240 8\,mi,
868 vs 728\,ft) or the same (ring\_300 3\,mi, 231\,ft).

\paragraph{Manual checklists.} \todo{Attach the completed US8 checklists, GPX
import record and route review from Appendix~\ref{app:checklists}.}
